\exercisesection{}

\begin{enumerate}
\def\labelenumi{\arabic{enumi})}
\item
  The function \(\varphi_0\) on \(\mathbf{R}\) defined by \(\varphi_0(x) = 0\) if \(x \leqslant 0\) and \(\varphi_0(x) = e^{-1/x}\) if \(x > 0\) is of class \(\mathrm{C}^\infty\) and satisfies \(\varphi_0^{(k)}(0) = 0\) for every integer \(k \in \mathbf{N}\) .
\item
  Let \(x_0 \in \mathbf{R}^n\) , \(r > 0\) and \(t > 1\) . There exists \(\varphi \in \mathcal{D}(\mathbf{R}^n)\) such that
\end{enumerate}

\begin{enumerate}
\def\labelenumi{(\roman{enumi})}
\item
  One has \(0 \leqslant \varphi \leqslant 1\) ;
\item
  One has \(\varphi(x)=1\) for \(\|x-x_{0}\|\leqslant r\) and \(\varphi(x)=0\) for \(\|x-x_{0}\|\geqslant tr\) ;
\item
  For every \(\alpha \in \mathbf{N}^n\) , we have
\end{enumerate}

\[
\sup _ {x \in \mathbf {R} ^ {n}} | \partial^ {\alpha} \varphi (x) | \leqslant \alpha ! ((t - 1) r) ^ {| \alpha |}.
\]

\begin{enumerate}
\def\labelenumi{\arabic{enumi})}
\setcounter{enumi}{2}
\tightlist
\item
  Let \(n \in \mathbf{N}\) . We denote by \(\mu_n\) the Lebesgue measure on \(\mathbf{R}^n\) .
\end{enumerate}

\begin{enumerate}
\def\labelenumi{\alph{enumi})}
\tightlist
\item
  Let \(\mathcal{H}_1\) be the image of the map \(\partial_1\colon \mathcal{D}(\mathbf{R}^n)\to \mathcal{D}(\mathbf{R}^n)\) ; it is the space of \(\varphi \in \mathcal{D}(\mathbf{R}^n)\) such that
\end{enumerate}

\[
\int_ {\mathbf {R}} \varphi (x, y) d \mu_ {1} (x) = 0
\]

for all \(y \in \mathbf{R}^{n-1}\) .

\begin{enumerate}
\def\labelenumi{\alph{enumi})}
\setcounter{enumi}{1}
\tightlist
\item
  Let \(\varphi_0\in \mathcal{D}(\mathbf{R})\) be such that its integral over \(\mathbf{R}\) is equal to 1. For every \(\varphi \in \mathcal{D}(\mathbf{R}^n)\) , there exist \(\varphi_{1}\in \mathcal{H}_{1}\) such that
\end{enumerate}

\[
\varphi (x, y) = \varphi_ {1} (x, y) + \varphi_ {0} (x) \int_ {\mathbf {R}} \varphi (t, y) d \mu_ {1} (t)
\]

for all \((x,y)\in \mathbf{R}\times \mathbf{R}^{n - 1}\).

\begin{enumerate}
\def\labelenumi{\alph{enumi})}
\setcounter{enumi}{2}
\tightlist
\item
  For every distribution \(f \in \mathcal{D}'(\mathbf{R}^{n})\) , there exists \(g \in \mathcal{D}'(\mathbf{R}^{n})\) such that \(\partial_{1}g = f\) . If \(g_{1}\) and \(g_{2}\) are in \(\mathcal{D}'(\mathbf{R}^{n})\) satisfy \(\partial_{1}g_{1} = \partial_{1}g_{2} = f\) , then there exists a distribution \(h \in \mathcal{D}'(\mathbf{R}^{n-1})\) such that
\end{enumerate}

\[
\langle g _ {1} - g _ {2}, \varphi \rangle = \langle h, \psi \rangle
\]

for all \(\varphi\in\mathcal{D}(\mathbf{R}^{n})\) , where \(\psi\in\mathcal{D}(\mathbf{R}^{n-1})\) is defined by

\[
\psi (y) = \int_ {\mathbf {R}} \varphi (t, y) d \mu_ {1} (t).
\]

\(d)\) Let U be a connected open subset of \(\mathbf{R}^{n}\). Let \(f \in \mathcal{D}'(\mathrm{U})\) be a distribution such that \(\partial_{i}f = 0\) for \(1 \leqslant i \leqslant n\) . The distribution \(f\) is the distribution associated with a constant function.

\begin{enumerate}
\def\labelenumi{\arabic{enumi})}
\setcounter{enumi}{3}
\tightlist
\item
  We denote by \(\mu\) the Lebesgue measure on \(\mathbf{R}\) .
\end{enumerate}

\begin{enumerate}
\def\labelenumi{\alph{enumi})}
\tightlist
\item
  Let \(\nu\) be a measure on \(\mathbf{R}\) and \(a \in \mathbf{R}\) such that \(\nu(\{a\}) = 0\) . For \(x \in \mathbf{R}\) we set
\end{enumerate}

\[
f (x) = \left\{ \begin{array}{l l} \nu ([ a, x ]) & \text { if } x \geqslant a, \\ \nu ([ x, a ]) & \text { if } x \leqslant a. \end{array} \right.
\]

The function f is continuous and the associated distribution satisfies \(f' = \nu\) .

\begin{enumerate}
\def\labelenumi{\alph{enumi})}
\setcounter{enumi}{1}
\item
  Let f be a distribution on R such that \(f'\) is a measure \(\nu\) . The distribution f is associated with the continuous function on R defined as in the preceding question.
\item
  Let \(f\) be a distribution on \(\mathbf{R}\) such that \(f^{(k)}\) is a measure for every integer \(k \geqslant 0\) . Then \(f\) is the distribution associated with a function in \(\mathcal{C}^{\infty}(\mathbf{R})\) .
\end{enumerate}

\begin{enumerate}
\def\labelenumi{\arabic{enumi})}
\setcounter{enumi}{4}
\item
  The family of distributions associated with point measures \(\varepsilon_{n}\) for \(n\in\mathbf{Z}\) is summable in \(\mathcal{S}^{\prime}(\mathbf{R})\). Its sum \(\omega\) satisfies \(\mathcal{F}(\omega)=\omega\) . (This formula is equivalent to the Poisson formula for the locally compact group \(\mathrm{G}=\mathbf{R}^{n}\) equipped with Lebesgue measure and the subgroup \(\mathrm{H}=\mathbf{Z}^{n}\) , cf. II, § 1, no. 7.)
\item
  Let A be a commutative unital algebra over C such that the following conditions are satisfied:
\end{enumerate}

\begin{enumerate}
\def\labelenumi{(\roman{enumi})}
\item
  The algebra \(\mathcal{C}(\mathbf{R})\) is a unital subalgebra of A;
\item
  There exists a derivation \(d\colon\mathrm{A}\to\mathrm{A}\) whose restriction to \(\mathcal{C}^{1}(\mathbf{R})\) coincides with the usual derivation of functions.
\end{enumerate}

\begin{enumerate}
\def\labelenumi{\alph{enumi})}
\item
  The identity map X of R is invertible in A. (Prove that the map \(g: \mathbf{R} \to \mathbf{R}\) such that \(g(0) = 0\) and \(g(x) = x(\log(|x|) - 1)\) for \(x \neq 0\) is continuous and satisfies \(gX = 1\) .)
\item
  One has \(d \circ d(|\mathrm{X}|) = 0\) .
\item
  There does not exist a bilinear map \(m\colon\mathcal{D}^{\prime}(\mathbf{R})\times\mathcal{D}^{\prime}(\mathbf{R})\to\mathcal{D}^{\prime}(\mathbf{R})\) such that \(m(f,g)=fg\) for every pair \((f,g)\in\mathcal{D}(\mathbf{R})^{2}\) , and such that
\end{enumerate}

\[
m (f, g) ^ {\prime} = m \left(f ^ {\prime}, g\right) + m \left(f, g ^ {\prime}\right)
\]

for all \((f,g)\in\mathcal{D}^{\prime}(\mathbf{R})^{2}\) ("impossibility of multiplication of distributions").

\begin{enumerate}
\def\labelenumi{\alph{enumi})}
\setcounter{enumi}{3}
\tightlist
\item
  Give examples of sequences \((f_{n})\) and \((g_{n})\) of test functions on R such that \(f_{n} \to \varepsilon_{0}\) and \(g_{n} \to \varepsilon_{0}\) in \(\mathcal{D}'(\mathbf{R})\) , and such that \(f_{n}^{2}\) and \(g_{n}^{2}\) converge in \(\mathcal{D}'(\mathbf{R})\) to different limits.
\end{enumerate}

\begin{enumerate}
\def\labelenumi{\arabic{enumi})}
\setcounter{enumi}{6}
\tightlist
\item
  Let \(n \in N\) . We denote by \(\widetilde{\gamma}\) the linear representation of \(R^{n}\) on the space of functions from \(R^{n}\) to C defined by
\end{enumerate}

\[
\widetilde {\gamma} (h) f (x) = f (x + h)
\]

for all \(h \in \mathbf{R}^n\) and every function \(f: \mathbf{R}^n \to \mathbf{C}\) .

\begin{enumerate}
\def\labelenumi{\alph{enumi})}
\item
  For every \(h \in R^{n}\) , the map \(\widetilde{\gamma}(h)\) defines, by passing to subspaces, an automorphism of the locally convex space of \(\mathcal{D}(\mathbf{R}^{n})\) , which is still denoted \(\widetilde{\gamma}(h)\) .
\item
  We denote \(\pmb{\gamma}\) the automorphism of \(\mathcal{D}'(\mathbf{R}^n)\) transposed to \(\widetilde{\pmb{\gamma}}(-h)\) . If \(f \in \mathcal{D}(\mathbf{R}^n)\) , then \(\pmb{\gamma}(h)f = \widetilde{\pmb{\gamma}}(h)f\) for all \(h \in \mathbf{R}^n\) .
\item
  The map \(h \mapsto \gamma(h)\) is a linear representation of \(\mathbf{R}^{n}\) in \(\mathcal{D}'(\mathbf{R}^{n})\) .
\end{enumerate}

\(d)\) For every \(h \in \mathbf{R}^n\) and every \(\alpha \in \mathbf{N}^n\) , we have \(\partial^\alpha \circ \boldsymbol{\gamma}(h) = \boldsymbol{\gamma}(h) \circ \partial^\alpha\) .

\begin{enumerate}
\def\labelenumi{\alph{enumi})}
\setcounter{enumi}{4}
\tightlist
\item
  Let i be an integer such that \(1 \leqslant i \leqslant n\) and \(e_{i}\) the \(i\)-th vector of the canonical basis of \(R^{n}\) . For every distribution \(f \in \mathcal{D}'(\mathbf{R}^{n})\) , we have
\end{enumerate}

\[
\lim_{\substack{h\to 0\\ h\neq 0}}\frac{1}{h} (\boldsymbol {\gamma}(he_{i})f - f) = \partial_{i}f
\]

in \(\mathcal{D}'(\mathbf{R}^n)\) endowed with the weak topology.

\begin{enumerate}
\def\labelenumi{\arabic{enumi})}
\setcounter{enumi}{7}
\item
  Let \(n \in \mathbf{N}\) and let U be an open subset of \(R^{n}\) . A distribution \(f \in \mathcal{D}'(U)\) is said to be positive if \(\langle f, \varphi \rangle \geqslant 0\) for every \(\varphi \in \mathcal{D}(U)\) such that \(\varphi \geqslant 0\) . Prove that f is positive if and only if f is the distribution associated with a positive measure on U.
\item
  Let \(n \in \mathbf{N}\) and \(U\) be an open subset of \(\mathbf{R}^n\) . Let \(k \in \mathbf{N}\) and \(p \in ]1, +\infty[\) . The space \(W^{k,p}(U)\) is reflexive.
\item
  Let \(n \in \mathbf{N}\) and \(U\) be an open subset of \(\mathbf{R}^n\) . Let \(k \in \mathbf{N}\) and \(p \in [1, +\infty[\). Let \(V \subset U\) be a relatively compact open subset.
\end{enumerate}

Let \(\varphi \in \mathcal{D}(\mathbf{R}^n)\) be a positive function with support contained in the unit ball of \(\mathbf{R}^n\) and integral 1. For every \(\varepsilon > 0\) and \(x \in \mathbf{R}^n\) we set \(\varphi_{\varepsilon}(x) = \varepsilon^{-n}\varphi(x/\varepsilon)\). Let \(f \in \mathrm{W}^{k,p}(\mathrm{U})\) .

\begin{enumerate}
\def\labelenumi{\alph{enumi})}
\item
  The restriction \(f_{\mathrm{V}}\) of \(f\) to V belongs to \(\mathrm{W}^{k,p}(\mathrm{V})\) .
\item
  Let \(\widetilde{f}\) be the function on \(\mathbf{R}^n\) which coincides with \(f\) on U and vanishes outside U. When \(\varepsilon \to 0\), the restriction of \(\widetilde{f} * \varphi_{\varepsilon}\) to V belongs to \(\mathrm{W}^{k,p}(\mathrm{V})\) and converges to \(f_{\mathrm{V}}\) in \(\mathrm{W}^{k,p}(\mathrm{V})\) .
\end{enumerate}

\begin{enumerate}
\def\labelenumi{\arabic{enumi})}
\setcounter{enumi}{10}
\tightlist
\item
  Let \(n \in \mathbf{N}\) and let U be an open subset of \(\mathbf{R}^n\) . Let \(k \in \mathbf{N}\) and \(p \in [1, +\infty[\) . We denote by \(\mathrm{H}^{k,p}(\mathrm{U})\) the closure of \(\mathcal{C}^{\infty}(\mathrm{U})\) in \(\mathrm{W}^{k,p}(\mathrm{U})\) .
\end{enumerate}

For \(x \in \mathrm{U}\) , let \(\delta(x)\) be the distance from \(x\) to the boundary of \(\mathrm{U}\) in \(\mathbf{R}^n\); let \(\mathrm{U}_{-1} = \mathrm{U}_0 = \varnothing\) and

\[
\mathrm{U} _ {k} = \{x \in \mathrm{U} | \| x \| <   k, \delta (x) > k ^ {- 1} \}
\]

for every integer \(k \geqslant 1\) . Set \(\mathrm{V}_k = \mathrm{U}_{k+1} \cap (\mathbf{R}^n - \overline{\mathrm{U}}_{k-1})\) for \(k \geqslant 1\). The open sets \(\mathrm{V}_k\) cover U; let \((\psi_k)_{k \geqslant 1}\) be a partition of unity subordinate to the covering \((\mathrm{V}_k)_{k \geqslant 1}\) consisting of functions in \(\mathcal{D}(\mathrm{U})\) .

\begin{enumerate}
\def\labelenumi{\alph{enumi})}
\tightlist
\item
  For \(k \geqslant 1\) , the function
\end{enumerate}

\[
\varphi_{k} = \sum_{\substack{j\geqslant 1\\ \operatorname{Supp}(\psi_{j})\subset \mathrm{V}_{k}}}\psi_{j}
\]

belongs to \(\mathcal{D}(\mathrm{U})\) and has support contained in \(\mathrm{V}_k\) ; we have

\[
\sum_ {k \geqslant 1} \varphi_ {k} = 1
\]

on U.

\begin{enumerate}
\def\labelenumi{\alph{enumi})}
\setcounter{enumi}{1}
\tightlist
\item
  Let \(f \in \mathrm{W}^{k,p}(\mathrm{U})\) and \(\varepsilon > 0\) . Let \(\varphi\) and \(\varphi_{\varepsilon}\) for \(\varepsilon > 0\) functions as in Exercise 10. For every \(k \geqslant 1\), there exists \(\varepsilon_{k} > 0\) such that
\end{enumerate}

\[
\left\| \varphi_ {\varepsilon_ {k}} * (\psi_ {k} f) - \psi_ {k} f \right\| _ {k, p} <   \frac {\varepsilon}{2 ^ {k}}.
\]

\begin{enumerate}
\def\labelenumi{\alph{enumi})}
\setcounter{enumi}{2}
\tightlist
\item
  The function
\end{enumerate}

\[
f _ {\varepsilon} = \sum_ {k \geqslant 1} \varphi_ {\varepsilon_ {k}} * (\psi_ {k} f)
\]

belongs to \(\mathcal{C}^{\infty}(\mathrm{U})\) and \(\| f - f_{\varepsilon}\|_{k,p} < \varepsilon\) .

\(d)\). The space \(\mathrm{H}^{k,p}(\mathrm{U})\) is equal to \(\mathrm{W}^{k,p}(\mathrm{U})\) .

\begin{enumerate}
\def\labelenumi{\alph{enumi})}
\setcounter{enumi}{4}
\item
  The space \(\mathrm{H}^{k,p}(\mathrm{U})\) is identified with the completion of the space \(\mathcal{C}^{k}(\mathrm{U})\) for the norm \(f \mapsto \|f\|_{k,p}\) .
\item
  Let V be an open subset of \(R^{n}\) and let \(\varphi\colon U\to V\) be a diffeomorphism of class \(C^{k}\) . Assume that the partial derivatives of order \(\leqslant k\) of \(\varphi\) (resp. of \(\varphi^{-1}\) ) are bounded and uniformly continuous on U (resp. on V). The map \(f\mapsto f\circ\varphi\) is an isomorphism of topological vector spaces from \(W^{k,p}(V)\) to \(W^{k,p}(U)\) . (Use the result of the previous question.)
\end{enumerate}

\begin{enumerate}
\def\labelenumi{\arabic{enumi})}
\setcounter{enumi}{11}
\tightlist
\item
  \begin{enumerate}
  \def\labelenumii{\alph{enumii})}
  \tightlist
  \item
    For all integers \(n \in \mathbf{N}\) and \(k \in \mathbf{N}\) , and for every real number \(p \geqslant 1\) , we have \(\mathrm{W}^{k,p}(\mathbf{R}^n) = \mathrm{W}_0^{k,p}(\mathbf{R}^n)\) .
  \end{enumerate}
\end{enumerate}

\begin{enumerate}
\def\labelenumi{\alph{enumi})}
\setcounter{enumi}{1}
\tightlist
\item
  Let \(p \geqslant 1\) be a real number. Let \(U \subset R^{2}\) be the set of \((x, y) \in \mathbf{R}^{2}\) such that 0 \textless{} \textbar x\textbar{} \textless{} 1 and 0 \textless{} y \textless{} 1. The characteristic function \(\varphi\) of \(U \cap (\mathbf{R}_{+}^{*} \times \mathbf{R})\) belongs to \(W^{1,p}(U)\) ; if \(\varepsilon > 0\) is sufficiently small, there does not exist a function f of class \(C^{1}\) in the neighborhood of \(\overline{U}\) such that \(\|f - \varphi\|_{1,p} < \varepsilon\) .
\end{enumerate}

\begin{enumerate}
\def\labelenumi{\arabic{enumi})}
\setcounter{enumi}{12}
\tightlist
\item
  Let U be an open subset of \(\mathbf{R}^n\) and \(j\colon \mathrm{U}\to \mathbf{R}^n\) be the canonical inclusion. For every function \(f\colon \mathrm{U}\to \mathbf{C}\) , we denote by \(j_{!}f\) the zero extension of \(f\) to \(\mathbf{R}^n\) , that is, the function from \(\mathbf{R}^n\) to \(\mathbf{C}\) which coincides with \(f\) in U and is zero outside U.
\end{enumerate}

Let \(k \geqslant 0\) be an integer and \(p \geqslant 1\) a real number.

\begin{enumerate}
\def\labelenumi{\alph{enumi})}
\item
  Let \(f \in \mathrm{W}^{k,p}(\mathrm{U})\) . For every \(\alpha \in \mathbf{N}^n\) such that \(|\alpha| \leqslant k\) , we have \(\partial^\alpha(j_{!}f) = j_{!}(\partial^\alpha(f))\) in \(\mathcal{D}'(\mathbf{R}^n)\) .
\item
  The map \(f \mapsto j_{!}f\) induced by passage to subspaces, is an isometric linear map from \(\mathrm{W}_{0}^{k,p}(\mathrm{U})\) into \(\mathrm{W}^{k,p}(\mathbf{R}^{n})\) .
\end{enumerate}

\begin{enumerate}
\def\labelenumi{\arabic{enumi})}
\setcounter{enumi}{13}
\tightlist
\item
  Let U be an open subset of \(R^{n}\) and let \(k \geqslant 1\) be an integer. Let \(p \geqslant 1\) be a real number. Assume that the complement \(R^{n} - U\) is not negligible for Lebesgue measure. Let j denote the canonical inclusion of U into \(R^{n}\) .
\end{enumerate}

\begin{enumerate}
\def\labelenumi{\alph{enumi})}
\item
  There exists a relatively compact open rectangle B in \(R^{n}\) such that \(B \cap U\) and \(B \cap (\mathbf{R}^{n} - \mathbf{U})\) are not negligible.
\item
  Let \(\widetilde{f} \in \mathcal{D}(\mathbf{R}^n)\) which equals 1 in \(B \cap U\). Let \(f\) be the restriction of \(\widetilde{f}\) to U. Then \(f\) does not belong to \(\mathrm{W}_0^{k,p}(\mathrm{U})\) (Otherwise, the extension by zero \(j_{!}f\) would belong to \(\mathrm{W}^{k,p}(\mathbf{R}^n)\) ; the restriction to B of \(j_{!}f\) would be a distribution such that \(\partial_i(j_{!}f) = 0\) for \(1 \leqslant i \leqslant n\) , hence would be the distribution associated with a constant.)
\end{enumerate}

\begin{enumerate}
\def\labelenumi{\arabic{enumi})}
\setcounter{enumi}{14}
\tightlist
\item
  Let \(n \geqslant 1\) be an integer and let p be a real number such that \(1 \leqslant p < n\) .
\end{enumerate}

\begin{enumerate}
\def\labelenumi{\alph{enumi})}
\item
  Let \(\alpha\) be a real number such that \(0 < \alpha < n/p - 1\). Let \(g\) be a function on \(\mathbf{R}_{+}^{*}\) , infinitely differentiable, positive, such that \(g(x) = x^{-\alpha}\) for \(x\) sufficiently small and \(g(x) = 0\) for \(x\) sufficiently large; the function from \(\mathbf{R}^{n}\) to \(\mathbf{C}\) defined by \(x \mapsto g(\|x\|)\) belongs to \(\mathrm{W}^{1,p}(\mathbf{R}^{n})\) .
\item
  Let \(i \mapsto q_{i}\) be a bijection from N to \(Q^{n}\) and let f be the map from \(R^{n}\) to C defined by
\end{enumerate}

\[
f (x) = \sum_ {i \in \mathbf {N}} 2 ^ {- i} g (\| x - q _ {i} \|)
\]

for \(x \in \mathbf{R}^n\) . The series defining \(f\) converges in \(\mathrm{W}^{1,p}(\mathbf{R}^n)\) ; for every open set U of \(\mathbf{R}^n\) the restriction of \(f\) to U is not almost everywhere differentiable. \(c)\) There exists a function \(f \in \mathrm{W}^{1,n}(\mathbf{R}^n)\) which is continuous at no point. (Similar method with the function \(x \mapsto \log(|\log(|x|)|)\) .)

\begin{enumerate}
\def\labelenumi{\arabic{enumi})}
\setcounter{enumi}{15}
\tightlist
\item
  Let \(n \geqslant 1\) and \(k \geqslant 1\) be integers. Let \(p \geqslant 1\) be a real number. Let U be a bounded open subset of \(\mathbf{R}^n\) .
\end{enumerate}

\begin{enumerate}
\def\labelenumi{\alph{enumi})}
\item
  Suppose that \(n \neq kp\) . Let q be a real number such that \(1 \leqslant q \leqslant np/(n-kp)\) . The inclusion of \(\mathcal{D}(U)\) in \(L^{q}(U)\) admits a continuous extension \(i\) from \(W_{0}^{k,p}(U)\) into \(L^{q}(U)\) .
\item
  Suppose that \(1 \leqslant q \leqslant np/(n-kp)\) . The linear map \(i\) is compact ("Rellich-Kondrachov theorem").
\item
  Suppose that \(k \geqslant 1\) and that \(q = np/(n - p)\) . The linear map i is not compact.
\end{enumerate}

\begin{enumerate}
\def\labelenumi{\arabic{enumi})}
\setcounter{enumi}{16}
\tightlist
\item
  Let \(n \geqslant 1\) be an integer. We denote by \(\mu\) the Lebesgue measure on \(\mathbf{R}^n\) .
\end{enumerate}

\begin{enumerate}
\def\labelenumi{\alph{enumi})}
\tightlist
\item
  A bounded subset B of \(\mathrm{L}^{2}(\mathbf{R}^{n},\mu)\) is compact if and only if, for every \(\varepsilon>0\) , there exists a compact subset K of \(R^{n}\) such that
\end{enumerate}

\[
\int_ {\mathbf {R} ^ {n} - K} | f | ^ {2} + \int_ {\mathbf {R} ^ {n} - K} | \widehat {f} | ^ {2} <   \varepsilon
\]

for all \(f\in \mathrm{B}\)

\begin{enumerate}
\def\labelenumi{\alph{enumi})}
\setcounter{enumi}{1}
\tightlist
\item
  Let \(k \in N\) . Let K be a compact subset of \(R^{n}\) . We denote by \(\mathrm{H}_{k}^{2}(\mathrm{K})\) the closed subspace of the Sobolev space \(\mathrm{H}_{k}^{2}(\mathbf{R}^{n})\) consisting of functions vanishing outside the compact set K. The canonical injection \(\mathrm{H}_{k}^{2}(\mathrm{K}) \to \mathrm{H}_{k}^{2}(\mathbf{R}^{n})\) is compact.
\end{enumerate}

\begin{enumerate}
\def\labelenumi{\arabic{enumi})}
\setcounter{enumi}{17}
\tightlist
\item
  Let \(n \in \mathbf{N}\) and \(U\) be an open subset of \(\mathbf{R}^n\) . The goal of this exercise is to prove that \(\mathcal{D}'(U)\) is a bornological space.
\end{enumerate}

If K is a compact subset of U, we equip \(\mathcal{C}_{\mathrm{K}}^{\infty}(\mathrm{U})'\) with the topology of bounded convergence. We denote by \(\pi_{\mathrm{K}}\colon\mathcal{D}'(\mathrm{U})\to\mathcal{C}_{\mathrm{K}}^{\infty}(\mathrm{U})'\) the linear map which assigns to a distribution its restriction to \(\mathcal{C}_{\mathrm{K}}^{\infty}(\mathrm{U})\) ; it is continuous.

Let \(S \subset \mathcal{D}'(U)\) be a convex balanced bornivorous set.

\begin{enumerate}
\def\labelenumi{\alph{enumi})}
\item
  There exists a compact subset K of U such that \(\mathrm{Ker}(\pi_{\mathrm{K}}) \subset \mathrm{S}\) . (Otherwise, consider an increasing sequence \((\mathrm{K}_{m})_{m \in \mathbf{N}}\) of compact subsets of U whose interiors cover U and choose \(f_{m}\) in \(\mathrm{Ker}(\pi_{\mathrm{K}_{m}}) - \mathrm{S}\) ; then the set of elements \(mf_{m}\) for \(m \in N\) is a bounded subset not absorbed by S.)
\item
  The set T of distributions \(f \in \mathcal{D}'(\mathrm{U})\) such that \(f + \operatorname{Ker}(\pi_{\mathrm{K}}) \subset \mathrm{S}\) is a convex balanced bornivorous set.
\item
  The set \(\pi_{\mathrm{K}}(\mathrm{T})\) is a convex balanced bornivorous subset of \(\mathcal{C}_{\mathrm{K}}^{\infty}(\mathrm{U})'\) . (Let A be a bounded subset of \(\mathcal{C}_{\mathrm{K}}^{\infty}(\mathrm{U})'\) ; show that there exists a neighborhood W of 0 in \(\mathcal{D}(\mathrm{U})\) such that \(\mathrm{A} \subset (\mathrm{W} \cap \mathcal{C}_{\mathrm{K}}^{\infty}(\mathrm{U}))^{\circ}\) and deduce that A is contained in \(\pi_{\mathrm{K}}(\mathrm{W}^{\circ})\) , hence is absorbed by \(\pi_{\mathrm{K}}(\mathrm{T})\) .)
\end{enumerate}

\(d)\) Deduce that S is a neighborhood of 0, and conclude that the space \(\mathcal{D}'(U)\) is bornological. (Note that \(\pi_{\mathrm{K}}(\mathrm{T})\) is a neighborhood of 0, and observe that S contains \(\pi_{\mathrm{K}}^{-1}(\pi_{\mathrm{K}}(\mathrm{T}))\) .)

\begin{enumerate}
\def\labelenumi{\arabic{enumi})}
\setcounter{enumi}{18}
\tightlist
\item
  Let \(n \in \mathbf{N}\) and let \(\mathrm{U} \subset \mathbf{R}^n\) be an open set.
\end{enumerate}

\begin{enumerate}
\def\labelenumi{\alph{enumi})}
\tightlist
\item
  Let \(f \in \mathcal{D}'(\mathrm{U})\) . The union V of the open sets \(\mathrm{W} \subset \mathrm{U}\) such that the restriction of \(f\) to W is zero is open in U.
\end{enumerate}

The support of f is called the complement of V; it is a closed subset of U.

\begin{enumerate}
\def\labelenumi{\alph{enumi})}
\setcounter{enumi}{1}
\item
  Let \(f \in \mathcal{D}'(\mathrm{U})\) and let \(\varphi \in \mathcal{D}(\mathrm{U})\) whose support does not meet the support of f. We have \(\langle f, \varphi \rangle = 0\) .
\item
  Let \(f \in \mathcal{D}'(\mathrm{U})\) and \(\alpha \in N^{n}\) . The support of \(\partial^{\alpha}f\) is contained in the support of f.
\item
  If \(f \in \mathcal{D}'(\mathrm{U})\) is the distribution associated with a measure \(\nu\) , then the support of f coincides with the support of \(\nu\) .
\item
  Let \(f\) be a distribution with compact support in U. The linear form \(f\) extends to a continuous linear form on the space \(\mathcal{C}^{\infty}(\mathrm{U})\) .
\item
  The space of distributions with compact support in U is identified with the dual of the space \(\mathcal{C}^{\infty}(\mathrm{U})\) .
\end{enumerate}

\(g)\) Let \(f \in \mathcal{D}'(\mathbf{R}^n)\) be a distribution with compact support. We have \(f \in \mathcal{S}'(\mathbf{R}^n)\) . \(h)\) Let \(f \in \mathcal{D}'(\mathbf{R}^n)\) be a distribution with compact support. The Fourier transform of \(f\) belongs to \(\mathcal{C}^\infty(\mathbf{R}^n)\) ; it is a function of polynomial growth together with all its derivatives.

\begin{enumerate}
\def\labelenumi{\arabic{enumi})}
\setcounter{enumi}{19}
\tightlist
\item
  Let \(n \in \mathbf{N}\) . We denote by \(V_n\) the open set in \(\mathbf{C}\) consisting of the complex numbers \(\lambda\) such that \(\mathcal{R}(\lambda) > -n - 1\) and \(\lambda \notin \{-1, -2, \ldots, -n\}\) .
\end{enumerate}

Let \(\lambda\in V_{0}\) . We denote by \(x_{+}^{\lambda}\) the distribution on \(\mathbf{R}\) defined by the locally integrable function

\[
x \mapsto \left\{ \begin{array}{l l} 0 & \text { if } x \leqslant 0, \\ x ^ {\lambda} & \text { if } x > 0. \end{array} \right.
\]

\begin{enumerate}
\def\labelenumi{\alph{enumi})}
\item
  The map \(\lambda \mapsto x_{+}^{\lambda}\) of \(\mathrm{V}_0\) in \(\mathcal{D}'(\mathbf{R})\) is analytic.
\item
  Let \(n \in \mathbf{N}\) . For every \(\lambda \in V_n\) , the map defined on \(\mathcal{D}(\mathbf{R})\) by
\end{enumerate}

\[
\begin{array}{l} \varphi \mapsto \int_ {0} ^ {1} x ^ {\lambda} \Big (\varphi (x) - \varphi (0) - x \varphi^ {\prime} (0) - \dots - \frac {x ^ {n - 1}}{(n - 1) !} \varphi^ {(n - 1)} (0) \Big) d x \\ \qquad + \int_ {1} ^ {+ \infty} x ^ {\lambda} \varphi (x) d x + \sum_ {k = 1} ^ {n} \frac {\varphi^ {(k - 1)} (0)}{(k - 1) ! (\lambda + k)} \end{array}
\]

is a tempered distribution; it does not depend on the integer n such that \(\lambda \in V_{n}\) , and coincides with \(x_{+}^{\lambda}\) when \(\lambda \in V_{0}\) .

We shall denote by \(x_{+}^{\lambda}\) the distribution thus defined for every \(\lambda\) belonging to the union of the sets \(V_{n}\) for \(n \in N\) .

\begin{enumerate}
\def\labelenumi{\alph{enumi})}
\setcounter{enumi}{2}
\tightlist
\item
  Let \(n \in N\) and let \(\lambda \in C\) such that \(-n - 1 < \mathcal{R}(\lambda) < -n\). For every Schwartz function \(\varphi\) on R, we have
\end{enumerate}

\[
\langle x _ {+} ^ {\lambda}, \varphi \rangle = \int_ {0} ^ {+ \infty} x ^ {\lambda} \left(\varphi (x) - \varphi (0) - x \varphi^ {\prime} (0) - \dots - \frac {x ^ {n - 1}}{(n - 1) !} \varphi^ {(n - 1)} (0)\right) d x.
\]

The map from \(V_{n}\) in \(\mathcal{D}'(\mathrm{U})\) defined by \(\lambda\mapsto x_{+}^{\lambda}\) is analytic.

\(d)\) For every \(\lambda\) belonging to the union of the sets \(V_{m}\) , we have the formula

\[
(x _ {+} ^ {\lambda}) ^ {\prime} = \lambda x _ {+} ^ {\lambda - 1}.
\]

\begin{enumerate}
\def\labelenumi{\arabic{enumi})}
\setcounter{enumi}{20}
\tightlist
\item
  \begin{enumerate}
  \def\labelenumii{\alph{enumii})}
  \tightlist
  \item
    Let \(f \in \mathcal{S}(\mathbf{R})\) . The function g on R such that \(g(0) = 2f'(0)\) and \(g(x) = (f(x) - f(-x))/x\) for all \(x \neq 0\) is integrable on R with respect to Lebesgue measure. The map that assigns to f the integral of g over R is a tempered distribution, denoted vp, and called the Cauchy principal value of 1/x.
  \end{enumerate}
\end{enumerate}

\begin{enumerate}
\def\labelenumi{\alph{enumi})}
\setcounter{enumi}{1}
\tightlist
\item
  For every \(f \in \mathcal{S}(\mathbf{R})\) , we have
\end{enumerate}

\[
\langle \mathrm{vp}, f \rangle = \lim _ {\varepsilon \rightarrow 0} \int_ {\mathbf {R} - [ - \varepsilon , \varepsilon ]} \frac {f (x)}{x} d x.
\]

\begin{enumerate}
\def\labelenumi{\alph{enumi})}
\setcounter{enumi}{2}
\item
  The restriction of vp to \(\mathbf{R} - \{0\}\) is the distribution associated with the locally integrable function \(x \mapsto 1/x\) on \(\mathbf{R} - \{0\}\) .
\item
  The distribution vp coincides with the derivative of the distribution associated with the locally integrable function on R taking value 0 at 0 and defined for \(x \neq 0\) by \(x \mapsto \log(|x|)\) .
\item
  Let \(\varepsilon > 0\) . We denote by \(p_{\varepsilon}\) the function \(x \mapsto 1 / (x + i\varepsilon)\) on \(\mathbf{R}\) . We have
\end{enumerate}

\[
\lim _ {\varepsilon \to 0} p _ {\varepsilon} = \mathrm{vp} - i \pi \varepsilon_ {0}
\]

in the space \(\mathcal{S}'(\mathbf{R})\) endowed with the weak topology, where \(\varepsilon_0\) is the tempered distribution associated with the point mass at 0.

\begin{enumerate}
\def\labelenumi{\arabic{enumi})}
\setcounter{enumi}{21}
\item
  The distribution on R associated with the exponential function is not tempered.
\item
  Let \(n \in \mathbf{N}\) and let U be an open set of \(\mathbf{R}^n\) . We denote by \(\mu\) the Lebesgue measure on U. For \(m \in \mathbf{N}\) , we denote \(\boldsymbol{m} = (m, \ldots, m) \in \mathbf{N}^n\) .
\end{enumerate}

\begin{enumerate}
\def\labelenumi{\alph{enumi})}
\tightlist
\item
  Let \(f \in \mathcal{D}(\mathrm{U})\) . Suppose that there exists an integer \(m \geqslant 1\) and a real number \(c \geqslant 0\) such that
\end{enumerate}

\[
| \langle f, \varphi \rangle | \leqslant c \int_ {\mathrm{U}} | \partial^ {\boldsymbol {m}} \varphi | d \mu
\]

for all \(\varphi\in\mathcal{D}(\mathrm{U})\). There then exists a function \(g\in\mathrm{L}^{\infty}(\mathrm{U})\) such that \(f=\partial^{m}g\) . (The linear map \(\varphi\mapsto\partial^{m}\varphi\) is injective on \(\mathcal{D}(\mathrm{U})\) ; let E be its image; consider the linear form on E into C that maps \(\partial^{m}\varphi\) to \(\langle f,\varphi\rangle\) and apply the Hahn--Banach theorem.)

\begin{enumerate}
\def\labelenumi{\alph{enumi})}
\setcounter{enumi}{1}
\tightlist
\item
  Let V be a relatively compact open set contained in U. There exists an integer \(m \geqslant 1\) and \(c \geqslant 0\) such that
\end{enumerate}

\[
| \langle f, \varphi \rangle | \leqslant c \int_ {\mathrm{U}} | \partial^ {\boldsymbol {m}} \varphi | d \mu
\]

for all \(\varphi\in\mathcal{D}(\mathrm{U})\) with support contained in V.

\begin{enumerate}
\def\labelenumi{\alph{enumi})}
\setcounter{enumi}{2}
\tightlist
\item
  For every relatively compact open set \(V \subset U\) , there exists a continuous function \(g \in \mathcal{C}(V)\) and \(\alpha \in N^{n}\) such that the restriction of f to V is equal to \(\partial^{\alpha}g\). \(d)\) Give an example of a distribution \(f \in \mathcal{D}'(\mathbf{R}^{n})\) such that f is not of the form \(f = \partial^{\alpha}g\) for a continuous function \(g \in \mathcal{C}(\mathbf{R}^{n})\) and \(\alpha \in N^{n}\) .
\end{enumerate}

\begin{enumerate}
\def\labelenumi{\arabic{enumi})}
\setcounter{enumi}{23}
\tightlist
\item
  Let \(n \in N\) and let U be an open set of \(R^{n}\). Let \(f \in \mathcal{D}'(U)\) be a distribution whose support C is compact in U.
\end{enumerate}

\begin{enumerate}
\def\labelenumi{\alph{enumi})}
\tightlist
\item
  There exists an integer \(m \in N\) such that the distribution f extends to a continuous linear form on the space \(\mathcal{C}^{m}(U)\) endowed with the topology of uniform convergence on compact subsets of functions and their partial derivatives of order \(\leqslant m\) .
\end{enumerate}

The smallest integer m satisfying this property is called the order of f.

\begin{enumerate}
\def\labelenumi{\alph{enumi})}
\setcounter{enumi}{1}
\item
  Let \(\varphi \in \mathcal{D}(\mathrm{U})\) such that \(\partial^{\alpha}\varphi\) vanishes on C for every \(\alpha \in \mathbf{N}^{n}\) such that \(|\alpha| \leqslant m\) . We then have \(\langle f, \varphi \rangle = 0\) . (Consider, for \(\delta > 0\) sufficiently small, the neighborhood \(\mathrm{U}_{\delta}\) of C consisting of the \(x \in \mathrm{U}\) such that \(d(x, \mathrm{C}) \leqslant \delta\) , and show that there exists \(\eta \geqslant 0\) , tending to 0 as \(\delta\) tends to 0, such that, for some \(\mathrm{C}_{0} \geqslant 0\) , one has \(|\partial^{\alpha}\varphi(x)| \leqslant \mathrm{C}_{0}\delta^{m-|\alpha|}\eta\) for \(x \in \mathrm{U}_{\delta}\) and \(|\alpha| \leqslant m\) ; show that there exists a function \(\beta_{\delta} \in \mathcal{D}(\mathrm{U})\) equal to 1 on \(\mathrm{U}_{\delta/4}\) and with support contained in \(\mathrm{U}_{\delta}\) satisfying \(|\partial^{\alpha}\beta_{\delta}| \leqslant \mathrm{C}_{1}\delta^{-|\alpha|}\) for some \(\mathrm{C}_{1} \geqslant 0\) ; finally observe that \(\langle f, \varphi \rangle = \langle f, \varphi\beta_{\delta} \rangle\) and that \(\varphi\beta_{\delta}\) tends to 0 in \(\mathcal{C}^{m}(\mathrm{U})\) .)
\item
  Let \(x \in U\) . The space of distributions f whose support of f is reduced to x has as an algebraic basis the distributions \(\partial^{\alpha}\varepsilon_{x}\) , where \(\alpha \in N^{n}\) .
\item
  Let \(n = 1\) . The formula
\end{enumerate}

\[
\langle f, \varphi \rangle = \lim _ {m \rightarrow + \infty} \left(\sum_ {1 \leqslant k \leqslant m} \varphi (k ^ {- 1}) - m \varphi (0) - \log (m) \varphi^ {\prime} (0)\right)
\]

defines a distribution on \(\mathbf{R}\) , of order 1, whose support is compact, equal to the set \(C = \{0\} \cup \{1/(n+1) \mid n \in \mathbf{N}\}\) . There exist sequences \((\varphi_n)\) in \(\mathcal{D}(\mathbf{R})\) such that \(\partial^{\alpha}\varphi_{n}\) converges uniformly to 0 on C for every \(\alpha\in\mathbf{N}^{n}\) such that \(|\alpha|\leqslant1\) , but \(\langle f,\varphi_{n}\rangle\to+\infty\) .

\begin{enumerate}
\def\labelenumi{\arabic{enumi})}
\setcounter{enumi}{24}
\tightlist
\item
  Let \(n \in \mathbf{N}\) . We denote by \(N\) the Euclidean norm on \(\mathbf{R}^n\) and let \(\mu\) be the Lebesgue measure.
\end{enumerate}

We denote by \(\mathcal{D}_{\mathrm{L}^{1}}(\mathbf{R}^{n})\) the vector subspace of \(\mathcal{C}^{\infty}(\mathbf{R}^{n})\) whose elements are the functions \(\varphi\) such that \(\partial^{\alpha}\varphi\in\mathrm{L}^{1}(\mathbf{R}^{n})\) and \(\partial^{\alpha}\varphi\in\mathcal{C}_{0}(\mathbf{R}^{n})\) for all \(\alpha\in\mathbf{N}^{n}\) . \(a)\) The space \(\mathcal{D}_{\mathrm{L}^{1}}(\mathbf{R}^{n})\) , endowed with the topology defined by the seminorms

\[
q _ {\alpha} (\varphi) = \int_ {\mathbf {R} ^ {n}} | \partial^ {\alpha} \varphi | d \mu , \quad \alpha \in \mathbf {N} ^ {n},
\]

is a Fréchet space.

\begin{enumerate}
\def\labelenumi{\alph{enumi})}
\setcounter{enumi}{1}
\item
  The space \(\mathcal{D}(\mathbf{R}^n)\) is contained in and is dense in \(\mathcal{D}_{\mathrm{L}^1}(\mathbf{R}^n)\) . One can identify the dual space \(\mathcal{D}_{\mathrm{L}^1}'(\mathbf{R}^n)\) of \(\mathcal{D}_{\mathrm{L}^1}(\mathbf{R}^n)\) with a subspace of \(\mathcal{D}'(\mathbf{R}^n)\) .
\item
  Let \(f \in \mathcal{D}'(\mathbf{R}^n)\) . We have \(f \in \mathcal{D}_{\mathrm{L}^1}'(\mathbf{R}^n)\) if and only if there exists a finite set I, a family \((\alpha_i)_{i \in \mathrm{I}}\) in \(\mathbf{N}^n\) and a family \((g_i)_{i \in \mathrm{I}}\) in \(\mathrm{L}^\infty(\mathbf{R}^n)\) such that
\end{enumerate}

\[
f = \sum_ {i \in I} \partial^ {\alpha_ {i}} g _ {i}.
\]

(Determine I and the family \((\alpha_{i})\) so that \(|\langle f,\varphi\rangle |\) is bounded when \(\varphi\) is such that \(q_{\alpha_i}(\varphi)\leqslant 1\) for all \(i\in \mathrm{I}\) ; identify \(\mathcal{D}_{\mathrm{L}^1}(\mathbf{R}^n)\) with a vector subspace of \(\mathrm{L}^1 (\mathbf{R}^n)^m\) by the mapping \(\varphi \mapsto (\partial^{\alpha_i}\varphi)_{i\in \mathrm{I}}\) , and apply the Hahn-Banach theorem.)

\begin{enumerate}
\def\labelenumi{\alph{enumi})}
\setcounter{enumi}{3}
\tightlist
\item
  Let \(f \in \mathcal{D}'(\mathbf{R}^n)\) . We have \(f \in \mathcal{D}_{\mathrm{L}^1}'(\mathbf{R}^n)\) if and only if there exists a continuous function \(g\) with moderate growth on \(\mathbf{R}^n\) and \(\alpha \in \mathbf{N}^n\) such that \(f = \partial^\alpha g\). \(e)\) Let \(f \in \mathcal{D}'(\mathbf{R}^n)\) . There exists an integer \(m \in \mathbf{N}\) such that \((1 + \mathrm{N}^2)^{-m/2}f\) belongs to \(\mathcal{D}_{\mathrm{L}^1}'(\mathbf{R}^n)\) .
\end{enumerate}

\(f)\) Let \(f \in \mathcal{D}'(\mathbf{R}^n)\) . We have \(f \in \mathcal{S}'(\mathbf{R}^n)\) if and only if there exists a bounded continuous function \(g\colon \mathbf{R}^n \to \mathbf{C}\) , an integer \(m \in \mathbf{N}\) and \(\alpha \in \mathbf{N}^n\) such that \(f = \partial^\alpha((1 + \mathrm{N}^2)^{m/2}g)\) .

\begin{enumerate}
\def\labelenumi{\arabic{enumi})}
\setcounter{enumi}{25}
\tightlist
\item
  Let \(n \in N\) . We equip \(R^{n}\) with Lebesgue measure \(\mu\) and identify \(R^{n}\) with its dual by the mapping \((x, y) \mapsto \exp(2i\pi x \cdot y)\) .
\end{enumerate}

\begin{enumerate}
\def\labelenumi{\alph{enumi})}
\item
  For any function \(f \in \mathcal{S}(\mathbf{R}^n)\) , the map \(g \mapsto f * g\) is an automorphism of the locally convex space of \(\mathcal{S}(\mathbf{R}^n)\) .
\item
  Let \(f \in \mathcal{S}'(\mathbf{R}^{n})\) and \(g \in \mathcal{S}(\mathbf{R}^{n})\) . The map \(\varphi \mapsto \langle f, \check{g} * \varphi \rangle\) is a tempered distribution, called the convolution of f and g. If \(f \in \mathcal{S}(\mathbf{R}^{n})\) , it coincides with the distribution associated with \(f * g\) . We shall still denote by \(f * g\) the tempered distribution thus defined.
\end{enumerate}

In the following questions, we fix \(f \in \mathcal{S}'(\mathbf{R}^n)\) and \(g \in \mathcal{S}(\mathbf{R}^n)\) . \(c)\) One has \(f * g \in \mathcal{C}^\infty(\mathbf{R})\) and

\[
\partial^ {\alpha} (f * g) = \partial^ {\alpha} (f) * g = f * \partial^ {\alpha} (g)
\]

for all \(\alpha\in\mathbf{N}^{n}\) .

\(d)\) For every \(y \in \mathbf{R}^n\), let \(g_y\) be the function \(x \mapsto g(y - x)\) on \(\mathbf{R}^n\) . The function \(f * g\) satisfies

\[
(f * g) (y) = \langle f, g _ {y} \rangle
\]

for all \(y \in R^{n}\) .

\begin{enumerate}
\def\labelenumi{\alph{enumi})}
\setcounter{enumi}{4}
\item
  Let \(h \in \mathcal{S}(\mathbf{R}^n)\) . We have \((f*g)*h = f*(g*h)\) .
\item
  The function \(f*g\) has polynomial growth, as do all its derivatives. (Use Exercise 25.)
\item
  One has \(\mathcal{F}(f*g) = \mathcal{F}(f)\mathcal{F}(g)\) .
\item
  The Fourier transform of fg is an infinitely differentiable function on \(R^{n}\) and has polynomial growth, as do all its derivatives. Moreover, for every \(y \in R^{n}\) , we have
\end{enumerate}

\[
\mathcal {F} (f g) (y) = \langle f, g e _ {- y} \rangle
\]

where \(e_y\) denotes the character \(x \mapsto \exp(2i\pi x \cdot y)\) of \(\mathbf{R}^n\) .

\begin{enumerate}
\def\labelenumi{\arabic{enumi})}
\setcounter{enumi}{26}
\tightlist
\item
  Let \(n \in \mathbf{N}\). We denote by N the Euclidean norm on \(\mathbf{R}^{n}\) and let \(\mu\) be the Lebesgue measure on every open subset of \(\mathbf{R}^{n}\). Let \(\Delta\) be the differential operator
\end{enumerate}

\[
\Delta = - \sum_ {i = 1} ^ {n} \partial_ {i} ^ {2}
\]

viewed as an endomorphism of \(\mathcal{D}'(\mathbf{R}^{n})\) .

For \(k \in \mathbf{R}\) , we denote by \(\mathrm{S}^k(\mathbf{R}^n)\) the space of tempered distributions \(f\) in \(\mathcal{S}'(\mathbf{R}^n)\) such that \((1 + \mathrm{N}^k)\mathcal{F}(f)\) belongs to \(\mathrm{L}^2(\mathbf{R}^n)\) .

Let \(U \subset R^{n}\) be an open subset of \(R^{n}\) . We denote by \(\mathrm{S}_{\mathrm{U}}^{k}(\mathbf{R}^{n})\) the space of distributions \(f \in \mathcal{D}'(\mathbf{R}^{n})\) such that for every test function \(\varphi \in \mathcal{D}(\mathbf{R}^{n})\) with support contained in U, we have \(\varphi f \in \mathrm{S}^{k}(\mathbf{R}^{n})\) .

\begin{enumerate}
\def\labelenumi{\alph{enumi})}
\tightlist
\item
  The space \(\mathrm{S}^{k}(\mathbf{R}^{n})\) is a Hilbert space with the inner product
\end{enumerate}

\[
\langle f _ {1} \mid f _ {2} \rangle = \int_ {\mathbf {R} ^ {n}} (1 + \mathrm{N} ^ {k}) \overline {{\mathcal {F} (f _ {1})}} \mathcal {F} (f _ {2}) d \mu .
\]

When \(k \in \mathbf{N}\) , we have \(W^{k}(\mathbf{R}^{n}) = S^{k}(\mathbf{R}^{n})\) and the norms of these two Banach spaces are equivalent; moreover, we then have \(W^{k}(\mathbf{R}^{n}) \subset S_{\mathrm{U}}^{k}(\mathbf{R}^{n})\) .

\begin{enumerate}
\def\labelenumi{\alph{enumi})}
\setcounter{enumi}{1}
\tightlist
\item
  The dual of \(\mathrm{S}^{k}(\mathbf{R}^{n})\) is identified with \(\mathrm{S}^{-k}(\mathbf{R}^{n})\) by the bilinear form
\end{enumerate}

\[
(f, g) \mapsto \int_ {\mathbf {R} ^ {n}} \mathcal {F} (f) (- x) \mathcal {F} (g) (x) d \mu (x)
\]

for \((f,g)\in \mathrm{S}^{-k}(\mathbf{R}^n)\times \mathrm{S}^k (\mathbf{R}^n).\)

\begin{enumerate}
\def\labelenumi{\alph{enumi})}
\setcounter{enumi}{2}
\tightlist
\item
  Suppose that U is relatively compact. We have
\end{enumerate}

\[
\bigcup_ {k \in \mathbf {R}} \mathrm{S} _ {\mathrm{U}} ^ {k} (\mathbf {R} ^ {n}) = \mathcal {D} ^ {\prime} (\mathbf {R} ^ {n}).
\]

(Let \(f \in \mathcal{D}'(\mathbf{R}^n)\) ; apply the results of Exercise 19 to \(\varphi f\) , where \(\varphi \in \mathcal{D}(\mathbf{R}^n)\) is equal to 1 on U.)

\(d)\) Let \(k \in \mathbf{R}\) and \(f \in \mathrm{S}^k(\mathbf{R}^n)\) . For every integer \(i\) such that \(1 \leqslant i \leqslant n\) , we have \(\partial_i f \in \mathrm{S}^{k-1}(\mathbf{R}^n)\) . Moreover, if \(\Delta(f) \in \mathrm{S}^k(\mathbf{R}^n)\) , then \(f \in \mathrm{S}^{k+2}(\mathbf{R}^n)\) .

\begin{enumerate}
\def\labelenumi{\alph{enumi})}
\setcounter{enumi}{4}
\tightlist
\item
  Let \(k \in \mathbf{R}\) and \(f \in \mathrm{S}_{\mathrm{U}}^{k}(\mathbf{R}^{n})\) . If \(\Delta(f) \in \mathrm{S}_{\mathrm{U}}^{k}(\mathbf{R}^{n})\) , then \(f \in \mathrm{S}_{\mathrm{U}}^{k+2}(\mathbf{R}^{n})\) .
\end{enumerate}

\(f)\) Let \(k > n/2\) be a real number and \(m \in \mathbf{N}\) be an integer such that \(m < k - n/2\) . The space \(\mathrm{S}^k(\mathbf{R}^n)\) is contained in \(\mathcal{C}^m(\mathbf{R}^n)\) (Using the Cauchy-Schwarz inequality, verify that \(x \mapsto x^\alpha \mathscr{F} f(x)\) belongs to \(\mathrm{L}^1(\mathbf{R}^n)\) for all \(\alpha \in \mathbf{N}^n\) such that \(|\alpha| \leqslant m\). Then show that \(\mathscr{F}(f)\) is differentiable on \(\mathbf{R}\) using the Fourier inversion formula.)

\(g)\) Let \(k > n/2\) be a real number and \(m \in \mathbf{N}\) be an integer such that \(m < k - n/2\). The restriction to U of an element of \(\mathrm{S}_{\mathrm{U}}^{k}(\mathbf{R}^{n})\) belongs to \(\mathcal{C}^{m}(\mathrm{U})\) (For every \(x \in \mathrm{U}\) , verify that one obtains a function \(g \in \mathcal{C}^{m}(\mathrm{U})\) well defined by setting \(g(x) = (\varphi_{x}f)(x)\) , where \(\varphi_{x} \in \mathcal{D}(\mathbf{R}^{n})\) is a function with support in U equal to 1 in a neighborhood of \(x\) ; finally verify that the restriction of \(f\) to U is equal to \(g\) .)

\begin{enumerate}
\def\labelenumi{\alph{enumi})}
\setcounter{enumi}{7}
\tightlist
\item
  Let \(g \in \mathcal{D}'(\mathbf{R}^n)\) . If \(f \in \mathcal{D}'(\mathbf{R}^n)\) satisfies \(\Delta(f) = g\) and if the restriction of \(g\) to U is in \(\mathcal{C}^\infty(\mathrm{U})\) , then the restriction of \(f\) to U is in \(\mathcal{C}^\infty(\mathrm{U})\) ("elliptic regularity theorem").
\end{enumerate}

\begin{enumerate}
\def\labelenumi{\arabic{enumi})}
\setcounter{enumi}{27}
\tightlist
\item
  Let \(n\) and \(m\) be integers, and let \(\mathrm{U} \subset \mathbf{R}^n\) and \(\mathrm{V} \subset \mathbf{R}^m\) be open sets. \(a)\) There exists a unique injective linear map \(u\) from \(\mathcal{D}(\mathrm{U}) \otimes \mathcal{D}(\mathrm{V})\) into \(\mathcal{D}(\mathrm{U} \times \mathrm{V})\) such that \(u(\varphi \otimes \psi)\) is the function \((x, y) \mapsto \varphi(x)\psi(y)\) for all \((\varphi, \psi) \in \mathcal{D}(\mathrm{U}) \times \mathcal{D}(\mathrm{V})\) . Moreover, the image of \(u\) is dense in \(\mathcal{D}(\mathrm{U} \times \mathrm{V})\) . One identifies the space \(\mathcal{D}(\mathrm{U}) \otimes \mathcal{D}(\mathrm{V})\) with a subspace of \(\mathcal{D}(\mathrm{U} \times \mathrm{V})\) by means of the map \(u\) .
\end{enumerate}

\begin{enumerate}
\def\labelenumi{\alph{enumi})}
\setcounter{enumi}{1}
\item
  For all distributions \(f\) on U and \(g\) on V, there exists a unique distribution \(f \otimes g\) on U × V such that \(\langle f \otimes g, \varphi \otimes \psi \rangle = \langle f, \varphi \rangle \langle g, \psi \rangle\) for all \((\varphi, \psi) \in \mathcal{D}(\mathrm{U}) \times \mathcal{D}(\mathrm{V})\) .
\item
  Let \(f \in \mathcal{D}(\mathrm{U})\) and \(g \in \mathcal{D}(\mathrm{V})\) . For every \(\eta \in \mathcal{D}(\mathrm{U} \times \mathrm{V})\) and every \(x \in U\), let \(\eta_x\) be the map \(y \mapsto \eta(x, y)\) . We have \(\eta_x \in \mathcal{D}(\mathrm{V})\) and the map \(\varphi\) which assigns to each \(x \in U\) \(\langle f, \eta_x \rangle\) belongs to \(\mathcal{D}(\mathrm{U})\) and satisfies
\end{enumerate}

\[
\langle f \otimes g, \eta \rangle = \langle f, \varphi \rangle .
\]

\begin{enumerate}
\def\labelenumi{\arabic{enumi})}
\setcounter{enumi}{28}
\tightlist
\item
  Let \(n\) and \(m\) be integers and let \(\mathrm{U} \subset \mathbf{R}^n\) and \(\mathrm{V} \subset \mathbf{R}^m\) be open sets. We resume the notation of the previous exercise. Let \(u\) be a continuous linear map from \(\mathcal{D}(\mathrm{U})\) into \(\mathcal{D}'(\mathrm{V})\) .
\end{enumerate}

\begin{enumerate}
\def\labelenumi{\alph{enumi})}
\item
  There exists a unique continuous linear form \(\widetilde{\kappa}\) defined on the subspace \(\mathcal{D}(\mathrm{U})\otimes \mathcal{D}(\mathrm{V})\) of \(\mathcal{D}(\mathrm{U}\times \mathrm{V})\) such that \(\widetilde{\kappa} (\varphi \otimes \psi) = \langle u(\varphi),\psi \rangle\) for all \((\varphi ,\psi)\) in \(\mathcal{D}(\mathrm{U})\times \mathcal{D}(\mathrm{V})\) .
\item
  There exists a distribution \(\kappa(u)\) on \(U \times V\) which extends \(\widetilde{\kappa}\) , and the map \(u \mapsto \kappa(u)\) from \(\mathcal{L}(\mathcal{D}(U), \mathcal{D}'(V))\) into \(\mathcal{D}'(U \times V)\) is continuous.
\item
  The map \(\kappa\) is bijective ("Schwartz kernel theorem"; to prove the surjectivity of \(\kappa\), given a distribution \(\eta\) on U × V, define \(u(\varphi)\) as the distribution \(\psi \mapsto \langle \eta, \varphi \otimes \psi \rangle\) on V.)
\end{enumerate}

\begin{enumerate}
\def\labelenumi{\arabic{enumi})}
\setcounter{enumi}{29}
\tightlist
\item
  Let \(n \in \mathbf{N}\) . We denote by \((x_1, \ldots, x_n)\) the dual basis of the canonical basis of \(\mathbf{R}^n\) .
\end{enumerate}

\begin{enumerate}
\def\labelenumi{\alph{enumi})}
\tightlist
\item
  Let \(c \geqslant 0\) and r \textgreater{} 0 be real numbers. There exists a function u defined and analytic on an open neighborhood U of 0 in \(R^{n}\) such that
\end{enumerate}

\[
\left(r - \sum_ {j = 1} ^ {n - 1} x _ {j} - u\right) \partial_ {n} u - c r \sum_ {j = 1} ^ {n - 1} \partial_ {j} u = c r
\]

on U and \(u(x,0)=0\) for \(x\in\mathbf{R}^{n-1}\) such that \((x,0)\in\mathrm{U}\) . (Consider first the case n=2.)

\begin{enumerate}
\def\labelenumi{\alph{enumi})}
\setcounter{enumi}{1}
\tightlist
\item
  Let U be an open neighborhood of 0 in \(\mathbf{R}^{n}\). Let \((a_{1},\ldots,a_{n-1},b)\) be functions from \(\mathrm{U}\times\mathbf{R}\) to \(\mathbf{C}\) that are analytic in a neighborhood of \((0,0)\in\mathrm{U}\times\mathbf{R}\). We denote by Q the map from the vector space of analytic functions on U to itself such that
\end{enumerate}

\[
\mathrm{Q} (u) = \partial_ {n} u - \sum_ {j = 1} ^ {n - 1} a _ {j} (x _ {1}, \dots , x _ {n - 1}, u) \partial_ {j} u + b (x _ {1}, \dots , x _ {n - 1}, u).
\]

There exists at most one function \(v\) defined and analytic in an open neighborhood \(V\) of 0 contained in \(U\) such that \(Q(v) = 0\) and such that \(v(x,0) = 0\) for all \(x \in \mathbf{R}^{n-1}\) such that \((x,0) \in V\) . (Verify that the coefficients of the Taylor expansion at 0 of such a solution \(v\) are uniquely determined.)

\begin{enumerate}
\def\labelenumi{\alph{enumi})}
\setcounter{enumi}{2}
\tightlist
\item
  Let \(f\) and \(g\) be formal power series in \(\mathbf{C}[[z_1, \ldots, z_n]]\) . One says that \(f\) majorizes \(g\) , and one writes \(f \ll g\) , if
\end{enumerate}

\[
f = \sum_ {\alpha \in \mathbf {N} ^ {n}} a _ {\alpha} z ^ {\alpha}, \qquad g = \sum_ {\alpha \in \mathbf {N} ^ {n}} b _ {\alpha} z ^ {\alpha}
\]

with \(|a_{\alpha}| \leqslant b_{\alpha}\) for all \(\alpha \in \mathbf{N}^{n}\) . If the formal series g converges in a neighborhood of 0 in \(R^{n}\) , then so does f.

\(d)\) If there exist functions \((\widetilde{a}_1, \ldots, \widetilde{a}_{n-1}, \widetilde{b})\) analytic in a neighborhood of 0 such that \(a_j \ll \widetilde{a}_j\) for \(1 \leqslant j \leqslant n-1\) and \(b \ll \widetilde{b}\) (in the sense that the corresponding Taylor expansions at 0 satisfy these conditions), and if the equation

\[
\partial_ {n} v - \sum_ {j = 1} ^ {n - 1} \widetilde {a} _ {j} (x _ {1}, \dots , x _ {n - 1}, v) \partial_ {j} v + \widetilde {b} (x _ {1}, \dots , x _ {n - 1}, v) = 0
\]

in the unknown \(v\) has an analytic solution in a neighborhood of 0, then the equation \(\mathrm{Q}(v) = 0\) has a unique analytic solution in a neighborhood of 0.

\begin{enumerate}
\def\labelenumi{\alph{enumi})}
\setcounter{enumi}{4}
\tightlist
\item
  There exists a unique solution of the equation \(Q(v) = 0\) that is analytic in a neighborhood of 0 (“Cauchy–Kowalewski theorem”).
\end{enumerate}

\begin{enumerate}
\def\labelenumi{\arabic{enumi})}
\setcounter{enumi}{30}
\tightlist
\item
  We denote by L the differential operator
\end{enumerate}

\[
\mathrm{L} = - i \partial_ {1} + \partial_ {2} - 2 (x _ {1} + i x _ {2}) \partial_ {3}
\]

of order 1 on \(\mathbf{R}^{3}\) , where \((x_{1}, x_{2}, x_{3})\) is the dual basis of the canonical basis of \(\mathbf{R}^{3}\), the \(x_{i}\) being viewed as polynomials on \(\mathbf{R}^{3}\) ("Lewy operator").

\begin{enumerate}
\def\labelenumi{\alph{enumi})}
\item
  Let U be an open neighborhood of 0 in \(\mathbf{R}^3\) . There exists a function \(\varphi \in \mathcal{D}(\mathrm{U})\) such that \(\varphi(0) = 1\) and such that \(\mathrm{L}(\varphi)\) is zero in a neighborhood V of 0 contained in U. (Use the Cauchy-Kowalewski theorem.)
\item
  Let \(u \in \mathcal{C}^{\infty}(\mathrm{U})\) be a function such that \(\mathrm{L}(u) = 0\) and \(u(0) = 0\). Assume that there exists \(\delta > 0\) such that \(\mathcal{R}(u(x)) < -2\delta\) for \(x \in U - V\). We set \(v_{t} = \varphi \exp(t(u + \delta))\). We have \(v_{t} \in \mathcal{C}^{\infty}(\mathrm{U})\) and \(\mathrm{L}(v_{t})\) tends to 0 in \(\mathcal{D}(\mathrm{U})\) as \(t \to +\infty\) .
\item
  Let F be the space of functions \(f \in \mathcal{C}^{\infty}(U)\) such that \(\partial^{\alpha}f\) is bounded on U for every \(\alpha \in N^{3}\). The space F, equipped with the seminorms \(q_{\alpha}(f) = \sup_{x \in U} |\partial^{\alpha}f|\) is a Fréchet space.
\end{enumerate}

\(d)\) Let \(\psi\in\mathcal{D}(\mathbf{R}^{3})\) . Denote by \(f_{t}\) the restriction to U of the function \(x\mapsto t^{3}\psi(tx)\) . We have \(f_{t}\in\mathrm{F}\) for \(t\geqslant1\) , and there exists a real number \(c\geqslant0\) such that \(q_{\alpha}(f_{t})\leqslant ct^{|\alpha|+3}\) for all \(\alpha\in\mathbf{N}^{3}\) and every \(t\geqslant1\) .

\begin{enumerate}
\def\labelenumi{\alph{enumi})}
\setcounter{enumi}{4}
\tightlist
\item
  For every \(t \geqslant 1\) , we denote by \(\lambda_{t}\) the linear form \(f \mapsto \int_{U} fv_{t}\) on F. It is continuous and for t sufficiently large, one has
\end{enumerate}

\[
\lambda_ {t} (f _ {t}) = e ^ {\delta t} \left(\int_ {\mathrm{U}} \psi (x) \exp \left(\sum_ {j = 1} ^ {3} \xi_ {j} x _ {j}\right) d x _ {1} d x _ {2} d x _ {3} + \varepsilon (t)\right)
\]

where \(\xi_{j} = \partial_{j}u(0)\) and \(\varepsilon (t)\to 0\) as \(t\rightarrow +\infty\).

\(f)\) Suppose that \((\xi_j) \neq 0\) ; there exists \(f \in \mathrm{F}\) such that the integrals \(\int_{\mathrm{U}} f v_t\) are not bounded when \(t \to +\infty\) . (If the assertion were false, obtain a contradiction by using the previous question and the Banach–Steinhaus theorem.)

\(g)\) The functions \(v_{t}\) do not converge to 0 in \(\mathcal{D}'(\mathrm{U})\) when \(t\to+\infty\) . \(h)\) Let \(\xi=(0,0,\xi_{3})\in\mathbf{R}^{3}\) such that \(\xi_{3}<0\). There exists a neighborhood U of 0 in \(\mathbf{R}^{3}\) and a function \(u\in\mathcal{C}^{\infty}(\mathrm{U})\) such that \(\mathrm{L}(u)=0\) and \(u(0)=0\) , which satisfies the conditions of questions \(b)\) and \(f)\) . (Use the Cauchy–Kowalewski theorem.)

\begin{enumerate}
\def\labelenumi{\roman{enumi})}
\tightlist
\item
  Let \((v_{t})\) be functions in \(\mathcal{D}(\mathrm{U})\) such that \(\mathrm{L}(v_{t})\to0\) , and \(f\in\mathcal{D}(\mathrm{U})\) such that \(\int_{\mathrm{U}}fv_{t}\) is not bounded when \(t\to+\infty\) (question g)); the equation Lv=f has no solution \(v\in\mathcal{D}'(\mathrm{U})\) ("Lewy's theorem"). (Note that the formal adjoint of L coincides with L.)
\end{enumerate}

\begin{enumerate}
\def\labelenumi{\arabic{enumi})}
\setcounter{enumi}{31}
\tightlist
\item
  Let \(n \geqslant 1\) be an integer. We denote by \(\Delta\) the Laplacian on \(\mathcal{D}'(\mathbf{R}^n)\), defined by
\end{enumerate}

\[
\Delta (f) = - \sum_ {i = 1} ^ {n} \partial_ {i} ^ {2} f.
\]

Let u be a locally integrable function on \(R^{n}\) such that \(\Delta(u)\) is a locally integrable function.

Let \(v\in \mathcal{L}^{\infty}(\mathbf{R}^{n})\) be the function defined by

\[
v (x) = \left\{ \begin{array}{l l} 0 & \text { if } u (x) = 0, \\ \overline {{u (x)}} / | u (x) | & \text { if } u (x) \neq 0. \end{array} \right.
\]

\begin{enumerate}
\def\labelenumi{\alph{enumi})}
\item
  For \(\varepsilon > 0\), set \(u_{\varepsilon} = (|u|^2 + \varepsilon^2)^{1/2}\) . Show that the distribution \(-\Delta(u_{\varepsilon}) - \mathcal{R}(\overline{u}/u_{\varepsilon} \Delta(u))\) is positive. (First consider the case where \(u \in \mathcal{C}^{\infty}(\mathbf{R}^{n})\) , then use convolution to treat the general case.)
\item
  The distribution \(-\Delta(|u|) - \mathcal{R}(v\Delta(u))\) is positive (“Kato inequality”).
\end{enumerate}

\begin{enumerate}
\def\labelenumi{\arabic{enumi})}
\setcounter{enumi}{32}
\tightlist
\item
  Let \(m \geqslant 1\) be an integer.
\end{enumerate}

\begin{enumerate}
\def\labelenumi{\alph{enumi})}
\tightlist
\item
  There exists a unique linear map \(A^{+}\) (resp. \(A^{-}\) ) from \(\mathrm{H}^{m}(\mathbf{R})\) to \(\mathrm{H}^{m-1}(\mathbf{R})\) such that
\end{enumerate}

\[
\mathrm{A} ^ {+} f (x) = f ^ {\prime} (x) - x f (x) \quad (\text { resp.   A } ^ {-} f (x) = f ^ {\prime} (x) + x f (x))
\]

for all \(f \in \mathcal{D}(\mathbf{R})\) and every \(x \in \mathbf{R}\) .

\begin{enumerate}
\def\labelenumi{\alph{enumi})}
\setcounter{enumi}{1}
\tightlist
\item
  The map \(A^{+}\) is a Fredholm operator of index 1 and \(A^{-}\) is a Fredholm operator of index -1. (Consider the images of the Hermite functions under \(A^{+}\) and \(A^{-}\), cf. Exercise 3 of II, p. 263.)
\end{enumerate}
% semantic-fragments: legacy-input-seam
\relax{}
